% PAGES: 220-221

\begin{examplex}
\footnote{A similar example for computing sample covariance and correlation
matrices for daily returns of Apple, Facebook, Google, Microsoft, and Yahoo was
included in Chapter 1. Details on computing percentage returns can also be found
there.}
The file \textit{data-DJ30-july2011-june2013.xlsx} from fepress.org/nla-primer
contains the end of week and end of month adjusted closing prices for eight
financial and technology Dow Jones components (AXP; BAC; JPM; CSCO; HPQ; IBM;
INTC; MSFT) between July 1, 2011 and June 30, 2013. We compute the sample
covariance and sample correlation matrix of the returns of the eight stocks using
this data.

The weekly (percentage) returns are computed using the end of week closing prices
for the eight stocks, which, using the notation from section 7.2, give the entries
of the matrix $T_{\bfX}$. After computing the mean normalized time series matrix
$\overline{T}_{\bfX}$ of the stock returns, the sample covariance matrix
$\widehat{\Sigma}_{\bfX,week}$ and the sample correlation matrix
$\widehat{\Omega}_{\bfX,week}$ are computed from the formulas (7.60) and (7.61),
respectively.
\[
\widehat{\Sigma}_{\bfX,week} \;=\;
\begin{pmatrix}
0.0010 & 0.0013 & 0.0010 & 0.0005 & 0.0007 & 0.0005 & 0.0006 & 0.0005 \\
0.0013 & 0.0036 & 0.0021 & 0.0010 & 0.0011 & 0.0009 & 0.0009 & 0.0009 \\
0.0010 & 0.0021 & 0.0019 & 0.0009 & 0.0011 & 0.0007 & 0.0007 & 0.0007 \\
0.0005 & 0.0010 & 0.0009 & 0.0015 & 0.0009 & 0.0006 & 0.0005 & 0.0005 \\
0.0007 & 0.0011 & 0.0011 & 0.0009 & 0.0034 & 0.0008 & 0.0010 & 0.0006 \\
0.0005 & 0.0009 & 0.0007 & 0.0006 & 0.0008 & 0.0009 & 0.0005 & 0.0004 \\
0.0006 & 0.0009 & 0.0007 & 0.0005 & 0.0010 & 0.0005 & 0.0011 & 0.0005 \\
0.0005 & 0.0009 & 0.0007 & 0.0005 & 0.0006 & 0.0004 & 0.0005 & 0.0009
\end{pmatrix}
\]
\[
\widehat{\Omega}_{\bfX,week} \;=\;
\begin{pmatrix}
1      & 0.6875 & 0.6995 & 0.3727 & 0.3607 & 0.4847 & 0.5286 & 0.5497 \\
0.6875 & 1      & 0.8068 & 0.4165 & 0.3253 & 0.4876 & 0.4517 & 0.4771 \\
0.6995 & 0.8068 & 1      & 0.5411 & 0.4300 & 0.5338 & 0.4848 & 0.5088 \\
0.3727 & 0.4165 & 0.5411 & 1      & 0.4002 & 0.5183 & 0.4222 & 0.4569 \\
0.3607 & 0.3253 & 0.4300 & 0.4002 & 1      & 0.4937 & 0.5135 & 0.3233 \\
0.4847 & 0.4876 & 0.5338 & 0.5183 & 0.4937 & 1      & 0.5482 & 0.4872 \\
0.5286 & 0.4517 & 0.4848 & 0.4222 & 0.5135 & 0.5482 & 1      & 0.5137 \\
0.5497 & 0.4771 & 0.5088 & 0.4569 & 0.3233 & 0.4872 & 0.5137 & 1
\end{pmatrix}
\]
Note that the rows and columns in the correlation and covariance matrices herein
correspond to the companies in the following order: AXP; BAC; JPM; CSCO; HPQ; IBM;
INTC; MSFT.

Similarly, the monthly returns are computed using the end of monthly closing
prices for the eight stocks. This gives the entries of the matrix $T_X$, which can
be used to compute the mean normalized time series matrix $\overline{T}_X$ of the
stock returns. The following sample covariance matrix and sample correlation
matrix are then computed by using formulas (7.60) and (7.61), respectively:
\[
\widehat{\Sigma}_{\bfX,month} \;=\;
\begin{pmatrix}
0.0030 & 0.0053 & 0.0047 & 0.0029 & 0.0039 & 0.0012 & 0.0017 & 0.0022 \\
0.0053 & 0.0201 & 0.0125 & 0.0058 & 0.0061 & 0.0022 & 0.0020 & 0.0047 \\
0.0047 & 0.0125 & 0.0124 & 0.0061 & 0.0049 & 0.0022 & 0.0023 & 0.0042 \\
0.0029 & 0.0058 & 0.0061 & 0.0076 & 0.0048 & 0.0020 & 0.0025 & 0.0028 \\
0.0039 & 0.0061 & 0.0049 & 0.0048 & 0.0193 & 0.0035 & 0.0039 & 0.0022 \\
0.0012 & 0.0022 & 0.0022 & 0.0020 & 0.0035 & 0.0020 & 0.0012 & 0.0008 \\
0.0017 & 0.0020 & 0.0023 & 0.0025 & 0.0039 & 0.0012 & 0.0039 & 0.0024 \\
0.0022 & 0.0047 & 0.0042 & 0.0028 & 0.0022 & 0.0008 & 0.0024 & 0.0034
\end{pmatrix}
\]
\[
\widehat{\Omega}_{\bfX,month} \;=\;
\begin{pmatrix}
1      & 0.6843 & 0.7659 & 0.5991 & 0.5098 & 0.4781 & 0.5116 & 0.6890 \\
0.6843 & 1      & 0.7923 & 0.4652 & 0.3100 & 0.3482 & 0.2269 & 0.5745 \\
0.7659 & 0.7923 & 1      & 0.6304 & 0.3188 & 0.4438 & 0.3388 & 0.6493 \\
0.5991 & 0.4652 & 0.6304 & 1      & 0.3947 & 0.5165 & 0.4610 & 0.5610 \\
0.5098 & 0.3100 & 0.3188 & 0.3947 & 1      & 0.5544 & 0.4464 & 0.2697 \\
0.4781 & 0.3482 & 0.4438 & 0.5165 & 0.5544 & 1      & 0.4360 & 0.3009 \\
0.5116 & 0.2269 & 0.3388 & 0.4610 & 0.4464 & 0.4360 & 1      & 0.6733 \\
0.6890 & 0.5745 & 0.6493 & 0.5610 & 0.2697 & 0.3009 & 0.6733 & 1
\end{pmatrix}
\]
As expected, the correlations between the returns of the financial companies, and
the correlations between the returns of the technology companies are generally
stronger than the correlations of the returns of a financial company and a
technology company. Also, all the returns, both weekly and monthly, are positively
correlated with each other.
\end{examplex}

An important issue in practice is that the sample covariance matrix of time series
data may be very close to singular, which affects the performance of many
algorithms, e.g., the Cholesky decomposition and linear solvers. The result of
Theorem 7.2 gives necessary and sufficient conditions for the sample covariance
matrix to be nonsingular, and will be used subsequently for the linear regression
of time series data; see section 8.2. This result is the time series version of
Lemma 7.5 for random variables.

\begin{theorem}
Let $X_1$, $X_2$, \ldots, $X_n$ be $n$ random variables given by the $N \times n$
time series data matrix $T_{\bfX} = \mathrm{col}(T_{X_k})_{k=1:n}$ at $N$ data
points, with $N > n$. The sample covariance matrix $\widehat{\Sigma}_{\bfX}$ is
nonsingular if and only if the vectors $T_{X_1}, T_{X_2}, \ldots, T_{X_n}, \bfone$
are linearly independent, i.e.,
\begin{equation}
\widehat{\Sigma}_{\bfX} \text{ nonsingular} \quad \Longleftrightarrow \quad T_{X_1},T_{X_2},\ldots,T_{X_n},\bfone \text{ linearly independent},
\end{equation}
where $\bfone$ denotes the $N \times 1$ column vector with all entries equal to $1$.
\end{theorem}

\begin{proof}
Recall from (7.60) that $\widehat{\Sigma}_{\bfX} = \frac{1}{N-1}\,\overline{T}^t_{\bfX}\overline{T}_{\bfX}$,
where $\overline{T}_{\bfX} = \mathrm{col}\left(\overline{T}_{X_k}\right)_{k=1:n}$ is
the mean--normalized $N \times n$ matrix of time series data for the $n$ random
variables; see (7.48) and (7.49).

Also, recall from Lemma 5.2 that a matrix of the form $M^tM$ is symmetric positive
definite if and only if the columns of $M$ are linearly independent.

Then, we conclude that
\begin{equation}
\widehat{\Sigma}_{\bfX} \text{ nonsingular} \quad \Longleftrightarrow \quad \overline{T}_{X_1},\overline{T}_{X_2},\ldots,\overline{T}_{X_n} \text{ linearly independent}.
\end{equation}
From (7.67), it follows that, in order to prove (7.66), it is enough to show that
\begin{align}
T_{X_1},T_{X_2},\ldots,T_{X_n},\bfone &\text{ linearly independent} \\
\Longleftrightarrow \quad \overline{T}_{X_1},\overline{T}_{X_2},\ldots,\overline{T}_{X_n} &\text{ linearly independent}.
\end{align}
\noindent$\bullet$ Assume that $T_{X_1}, T_{X_2}, \ldots, T_{X_n}, \bfone$ are linearly
independent. We will show that $\overline{T}_{X_1}, \overline{T}_{X_2}, \ldots,
\overline{T}_{X_n}$ are linearly independent, which is equivalent to showing that,
if $c_k$, $k=1:n$, are constants such that
\begin{equation}
\sum_{k=1}^n c_k\overline{T}_{X_k} \;=\; 0,
\end{equation}
then $c_k = 0$ for all $k = 1:n$.
% CONTINUES: \begin{proof} of Theorem 7.2 (opened above) and the surrounding bullet
% item are left OPEN here at the end of PDF page 221 (folio 205). The next file
% (sec07_c2_3.tex, PAGES 222-223) picks up with "Since $\overline{T}_{X_k} =
% (X_k(t_i)-\widehat{\mu}_{X_k})_{i=1:N}$ ..." and eventually closes this proof
% with \end{proof}. Do not open a new \begin{proof} there.
